\begin{frame}{``분산이 크다''가 실전에서 무엇을 의미하나}
  수학식으로는 ``분산 $\sigma^2$'' 한 줄이지만, 실전에서는
  이렇게 느껴진다:
  \begin{itemize}
    \item \kb{학습 곡선이 ``안 정해진다''}: 에피소드 1만 개
     까지 평균이 $0.85 \to 0.92 \to 0.87 \to 0.91$ 로
      출렁인다. ``아직 안 수렴한 건가, 이미 수렴한 건가''를
      구분하기 어렵다.
    \item \kb{한두 에피소드가 전체를 지배한다}: $\rho$ 가
      극단적으로 큰 에피소드 1$\sim$2개가 평균의 절반 이상을
      차지할 수 있다 (시뮬레이션에서 $\rho \cdot G_0$ 값이
      가장 큰 상위 1\% 에피소드가 전체 합의 상당 부분을
      담당). ``학습이 한두 운 좋은 게임의 결과에 좌우된다''.
    \item \kb{$\varepsilon$ 을 키우면 더 악화된다}:
      $\varepsilon$-greedy 의 $\varepsilon$ 이 크면 $b$ 와
      $\pi$ 의 차이가 커져 $\rho$ 의 분포가 더 넓어진다.
      ``탐험을 많이 하면 데이터가 좋아진다''는 on-policy
      직관이 off-policy 에서는 \textbf{탐험이 많을수록
      분산이 커진다}는 반대 방향으로 작용한다.
  \end{itemize}
\end{frame}
